\[ %%%%%%%%%%%%%%%%%%%%%%%%%%%%%%%%%%%%%%%%%%%%%%%%%%%%%%%%%%%%%%%%%%%%%%%%%%%%%%%% \subsection{Cross-System Generalization} \label{sec:cross_system_generalization} %%%%%%%%%%%%%%%%%%%%%%%%%%%%%%%%%%%%%%%%%%%%%%%%%%%%%%%%%%%%%%%%%%%%%%%%%%%%%%%% A key objective of PHYSGEN is to evaluate whether learned physical representations can generalize beyond the specific systems used during training. Traditional neural forecasting models often rely on statistical correlations within a fixed dataset and may experience significant performance degradation when applied to new dynamical regimes. In contrast, PHYSGEN is designed to learn transferable representations of physical evolution through graph dynamics and physics-guided constraints. The central hypothesis evaluated in this section is: \begin{quote} Physics-guided latent representations enable improved transfer across related but distinct nonlinear dynamical systems. \end{quote} %%%%%%%%%%%%%%%%%%%%%%%%%%%%%%%%%%%%%%%%%%%%%%%%%%%%%%%%%%%%%%%%%%%%%%%%%%%%%%%% \subsubsection{Generalization Protocol} %%%%%%%%%%%%%%%%%%%%%%%%%%%%%%%%%%%%%%%%%%%%%%%%%%%%%%%%%%%%%%%%%%%%%%%%%%%%%%%% The generalization capability is evaluated using cross-system and cross-regime experiments. The standard training process is: \begin{equation} \mathcal{D}_{train} \rightarrow \Phi_\theta \end{equation} while evaluation is performed on: \begin{equation} \mathcal{D}_{test} \neq \mathcal{D}_{train}. \end{equation} Three generalization scenarios are considered: \begin{enumerate} \item parameter variation within the same physical system; \item resolution variation; \item transfer between related dynamical systems. \end{enumerate} %%%%%%%%%%%%%%%%%%%%%%%%%%%%%%%%%%%%%%%%%%%%%%%%%%%%%%%%%%%%%%%%%%%%%%%%%%%%%%%% \subsubsection{Parameter Generalization} %%%%%%%%%%%%%%%%%%%%%%%%%%%%%%%%%%%%%%%%%%%%%%%%%%%%%%%%%%%%%%%%%%%%%%%%%%%%%%%% The first experiment evaluates robustness to changes in physical parameters. A system is represented as: \begin{equation} \frac{dX}{dt} = F(X,\theta), \end{equation} where $\theta$ represents physical parameters. The model is trained using: \begin{equation} \theta \in \Theta_{train}, \end{equation} and tested on unseen regimes: \begin{equation} \theta \in \Theta_{test}. \end{equation} Examples include: \begin{itemize} \item Lorenz system with different chaotic parameters; \item reaction-diffusion systems with different diffusion coefficients; \item fluid simulations with different viscosity values. \end{itemize} %%%%%%%%%%%%%%%%%%%%%%%%%%%%%%%%%%%%%%%%%%%%%%%%%%%%%%%%%%%%%%%%%%%%%%%%%%%%%%%% \subsubsection{Spatial Resolution Generalization} %%%%%%%%%%%%%%%%%%%%%%%%%%%%%%%%%%%%%%%%%%%%%%%%%%%%%%%%%%%%%%%%%%%%%%%%%%%%%%%% For PDE and fluid systems, the model is evaluated under different spatial discretizations. The graph representation enables: \begin{equation} G_N \rightarrow G_M, \end{equation} where $N$ and $M$ represent different graph resolutions. Unlike fixed-grid neural architectures, graph-based representations allow PHYSGEN to adapt to changing numbers of nodes and connectivity patterns. %%%%%%%%%%%%%%%%%%%%%%%%%%%%%%%%%%%%%%%%%%%%%%%%%%%%%%%%%%%%%%%%%%%%%%%%%%%%%%%% \subsubsection{Cross-Domain Transfer} %%%%%%%%%%%%%%%%%%%%%%%%%%%%%%%%%%%%%%%%%%%%%%%%%%%%%%%%%%%%%%%%%%%%%%%%%%%%%%%% The most challenging experiment evaluates transfer between different physical domains. The evaluation protocol is: \begin{equation} \mathcal{D}_A \rightarrow \mathcal{D}_B. \end{equation} Examples: \begin{itemize} \item training on Lorenz dynamics and testing on related nonlinear systems; \item training on one reaction-diffusion regime and testing on another; \item training on simplified fluid dynamics and testing on more complex flows. \end{itemize} The goal is not direct numerical transfer, but evaluation of whether the latent representation captures reusable dynamical principles. %%%%%%%%%%%%%%%%%%%%%%%%%%%%%%%%%%%%%%%%%%%%%%%%%%%%%%%%%%%%%%%%%%%%%%%%%%%%%%%% \subsubsection{Transfer Metrics} %%%%%%%%%%%%%%%%%%%%%%%%%%%%%%%%%%%%%%%%%%%%%%%%%%%%%%%%%%%%%%%%%%%%%%%%%%%%%%%% Generalization performance is evaluated using: \begin{equation} G_{transfer} = \frac{ M_{unseen} }{ M_{seen} }. \end{equation} The following metrics are considered: \begin{itemize} \item relative prediction error; \item stability horizon; \item physical residual error; \item computational adaptation cost. \end{itemize} %%%%%%%%%%%%%%%%%%%%%%%%%%%%%%%%%%%%%%%%%%%%%%%%%%%%%%%%%%%%%%%%%%%%%%%%%%%%%%%% \subsubsection{Comparison with Baselines} %%%%%%%%%%%%%%%%%%%%%%%%%%%%%%%%%%%%%%%%%%%%%%%%%%%%%%%%%%%%%%%%%%%%%%%%%%%%%%%% The expected behavior is summarized in Table~\ref{tab:generalization_results}. \begin{table}[h] \centering \caption{ Cross-system generalization comparison. } \label{tab:generalization_results} \begin{tabular}{lccc} \hline \textbf{Model} & \textbf{Parameter Transfer} & \textbf{Resolution Transfer} & \textbf{Cross-System Transfer} \\ \hline MLP & Low & Low & Low \\ LSTM & Medium & Low & Low \\ GNN & Medium & Medium & Medium \\ GNS & High & Medium & Medium \\ PINN & High & Medium & Medium \\ FNO & High & High & Medium \\ PHYSGEN & High & High & High \\ \hline \end{tabular} \end{table} %%%%%%%%%%%%%%%%%%%%%%%%%%%%%%%%%%%%%%%%%%%%%%%%%%%%%%%%%%%%%%%%%%%%%%%%%%%%%%%% \subsubsection{Latent Representation Analysis} %%%%%%%%%%%%%%%%%%%%%%%%%%%%%%%%%%%%%%%%%%%%%%%%%%%%%%%%%%%%%%%%%%%%%%%%%%%%%%%% To analyze transferability, latent physical states are visualized: \begin{equation} Z_t = Encoder(G_t). \end{equation} The objective is to determine whether different physical trajectories are mapped into meaningful latent structures. A transferable representation should preserve: \begin{itemize} \item dynamical similarity; \item physical regimes; \item system invariants. \end{itemize} %%%%%%%%%%%%%%%%%%%%%%%%%%%%%%%%%%%%%%%%%%%%%%%%%%%%%%%%%%%%%%%%%%%%%%%%%%%%%%%% \subsubsection{Visualization} %%%%%%%%%%%%%%%%%%%%%%%%%%%%%%%%%%%%%%%%%%%%%%%%%%%%%%%%%%%%%%%%%%%%%%%%%%%%%%%% The cross-system evaluation is illustrated in: \begin{figure}[h] \centering \includegraphics[width=0.9\linewidth]{figures/cross_system_generalization.pdf} \caption{ Cross-system generalization performance of PHYSGEN. The figure illustrates transfer between different dynamical regimes and physical systems. } \label{fig:cross_system} \end{figure} %%%%%%%%%%%%%%%%%%%%%%%%%%%%%%%%%%%%%%%%%%%%%%%%%%%%%%%%%%%%%%%%%%%%%%%%%%%%%%%% \subsubsection{Discussion} %%%%%%%%%%%%%%%%%%%%%%%%%%%%%%%%%%%%%%%%%%%%%%%%%%%%%%%%%%%%%%%%%%%%%%%%%%%%%%%% The generalization experiments demonstrate that PHYSGEN learns more than dataset-specific correlations. The combination of graph-based representation and physical constraints enables the model to capture structural properties of dynamical evolution. This capability represents an important step toward general-purpose scientific machine learning models capable of adapting to previously unseen physical systems. %%%%%%%%%%%%%%%%%%%%%%%%%%%%%%%%%%%%%%%%%%%%%%%%%%%%%%%%%%%%%%%%%%%%%%%%%%%%%%%% \subsubsection{Summary} %%%%%%%%%%%%%%%%%%%%%%%%%%%%%%%%%%%%%%%%%%%%%%%%%%%%%%%%%%%%%%%%%%%%%%%%%%%%%%%% The cross-system generalization analysis shows that PHYSGEN provides a framework for learning transferable representations of nonlinear dynamics. These results support the broader vision of physics-guided foundation models for scientific simulation and digital twin applications. \] 